\documentclass[11pt]{article}
\usepackage[utf8]{inputenc}
\usepackage[T1]{fontenc}
\usepackage{amsmath,amssymb}
\usepackage{graphicx}
\usepackage{booktabs}
\usepackage[margin=1in]{geometry}
\usepackage{hyperref}

\title{Sintesis: Parametro Alpha y Comparaciones con Tesis de Nico\\
\large (Notebook: combined\_findings)}
\author{}
\date{}

\begin{document}
\maketitle

\begin{abstract}
Este documento cubre dos cosas:
(1) el parametro $\alpha$ que controla como el solver hibrido
elige pools, y
(2) comparaciones detalladas con resultados de la tesis de Nico
sobre testing clasico.
Los arboles de decision basicos estan en el documento 01,
y el parametro $\beta$ y la verificacion de Gibbs en el documento 02.
\end{abstract}

%% ====================================================================
\section{El Parametro Alpha ($\alpha$): Aclarar vs Informar en el Solver Hibrido}

\subsection{Contexto: que es el solver hibrido?}

El solver hibrido hace $K$ pasos con greedy y despues cambia a DP exacto
(ver documento 01, seccion 10).
Durante los pasos greedy, necesita decidir que pool testear.

\subsection{Que es $\alpha$?}

$\alpha$ controla como el solver hibrido puntua cada pool candidato:

$$\text{score}_\alpha = \alpha \times \underbrace{P(r=0) \cdot \sum u_i}_{\text{aclarar}} \;+\; (1 - \alpha) \times \underbrace{\text{ganancia\_de\_info}}_{\text{aprender}}$$

\begin{itemize}
  \item $\alpha = 1$: solo le importa aclarar (puntaje miope estandar).
  \item $\alpha = 0$: solo le importa aprender (maximiza reduccion de incertidumbre).
  \item $\alpha = 0.5$: mitad y mitad.
\end{itemize}

\textbf{Analogia:} $\alpha$ es como elegir entre ``preguntar algo facil en un examen''
($\alpha = 1$, seguro sacas puntos) vs ``preguntar algo que te ayuda a entender
el resto del examen'' ($\alpha = 0$, inviertes en informacion).

\subsection{Diferencia entre $\alpha$ y $\beta$}

\begin{table}[htbp]
  \centering
  \begin{tabular}{lcc}
    \toprule
     & $\alpha$ (documento 03) & $\beta$ (documento 02) \\
    \midrule
    Donde se usa & Solver hibrido & Greedy puro \\
    Como funciona & Pondera aclarar vs info & Suma bonus de info al score \\
    Formula & $\alpha \cdot \text{aclarar} + (1-\alpha) \cdot \text{info}$ & $\text{aclarar} + \beta \cdot \text{info}$ \\
    Rango & $[0, 1]$ & $[0, \infty)$ \\
    Codigo & \texttt{hybrid\_solver.py} & \texttt{infection\_reward\_greedy.py} \\
    \bottomrule
  \end{tabular}
  \caption{$\alpha$ y $\beta$ son similares pero se usan en contextos distintos.}
\end{table}

\subsection{Nota: hay otro $\alpha$ en semi\_utility}

El archivo \texttt{semi\_utility.py} define un $\alpha$ diferente:
$$U_{\text{semi}}(\alpha) = \sum u_i \left[ \alpha \cdot P(\text{sano}_i \mid \text{historia}) + (1-\alpha) \cdot \mathbf{1}_{\text{aclarado}}(i) \right]$$

Aqui $\alpha = 0$ es utilidad por aclarar (binario: aclarado o no),
y $\alpha = 1$ es utilidad por probabilidad de estar sano.
\textbf{No confundir} con el $\alpha$ del solver hibrido.
En este documento nos referimos solo al $\alpha$ del solver hibrido.

\subsection{Experimento: EU vs $\alpha$}

$n=6$, $B=4$, $G=3$, $p=[0.3]^6$, $u=[1]^6$.
Probamos $\alpha \in \{0.0, 0.25, 0.5, 0.75, 1.0\}$.

\textbf{Resultado:} el mejor $\alpha$ depende de la instancia.
Generalmente, un valor intermedio ($\alpha \approx 0.5$) funciona bien.
$\alpha = 0$ (solo informacion) y $\alpha = 1$ (solo aclarar)
suelen ser peores que valores intermedios.

\subsection{Resumen de $\alpha$}

$\alpha$ intermedio ($\sim 0.5$) funciona bien en general.
A prevalencia extrema (muy baja o muy alta), no importa mucho.
Para detalles de $\beta$, ver documento 02.

%% ====================================================================
\section{Comparaciones con la Tesis de Nico}

La tesis de Nico estudia pooled testing clasico (resultado binario).
Nosotros extendemos al caso aumentado (resultado = conteo).
Reproducimos sus ejemplos clave para ver que cambia.

\subsection{Ejemplo 4.4: Greedy elige el pool equivocado}

$n=3$, $B=2$, $G=3$, $p = [1/4, 1/4, 1/4]$, $u = [1, 1, 1]$.

Nico mostro que greedy clasico puede elegir un mal pool de seguimiento
despues de un resultado positivo.

En nuestro framework aumentado, el optimo dinamico y el clasico dan
la misma EU ($= 1.875$) en este caso --- el espacio es tan chico ($n=3$)
que no hay diferencia.
Sin embargo, \textbf{greedy secuencial} ($= 1.6875$) si pierde calidad,
confirmando que la historia completa importa.

\textbf{Resultado:} optimo $>$ greedy; aumentado $=$ clasico en este caso chico.

\subsection{Ejemplo 6.2: Trade-off entre EU media y riesgo}

$n=3$, $B=2$, $G=3$, $p = [0.9, 0.5, 0.5]$, $u = [1, 1, 1]$.

Nota: $p_i$ es probabilidad de infeccion, asi que persona A ($p=0.9$)
esta \emph{casi seguro infectada}.

En este caso todas las estrategias dan $\text{EU} = 1.0$.
El espacio es demasiado chico para que haya diferencias.
El punto de Nico (trade-off entre EU media y riesgo de cola)
se observa mejor a $n$ mas grande, como en nuestro Part 6.

\subsection{Ejemplo 6.3: Efecto de G}

$n=4$, $B=2$, $p = [1.0, 0.75, 0.75, 0.75]$, $u = [1, 1, 1, 1]$.
Persona A esta seguro infectada ($p=1.0$). Variamos $G$.

En este caso extremo ($p_0 = 1$), las estrategias coinciden
($\text{EU} = 0.5$ para todo $G$).
El conteo aumentado no agrega informacion cuando una persona
ya es seguro infectada.

\subsection{Escalamiento con B}

$n=5$, $G=3$, $B = 1, 2, 3, 4$. 30 instancias aleatorias por $B$.

La ventaja del aumentado sobre el clasico es mayor a $B$ intermedio.
A $B$ grande (cercano a $n$), todos los individuos se resuelven
y las estrategias convergen.
Este resultado se confirma con mas detalle en el documento principal
(paper\_findings, Part 3B).

\subsection{Ejemplo 6.1: maxima divergencia}

$n=3$, $B=2$, $G=3$, $p = [0.91, 0.17, 0.69]$, $u = [0.5, 0.9, 0.33]$.

Nico identifica esta instancia como la de mayor gap dinamico vs estatico.
En nuestro framework, todas las estrategias dan EU similar ($\approx 0.849$)
en este caso chico.
La divergencia que Nico observa se da en el framework clasico con
diferentes restricciones de pool.

\subsection{El problema de la ``persona segura''}

La estrategia dinamica tiende a incluir a la persona menos infectada
como ancla en multiples pools.
Funciona bien en promedio, pero si esa persona esta infectada,
todos los pools que la incluian se invalidan.

Los datos muestran que con $\beta = 0$ o $\beta = 0.5$ el algoritmo
ya excluye al ancla en algunos casos.
Con $\beta \geq 1$, el algoritmo puede incluirlo mas (no menos),
lo cual baja la EU.
El efecto de $\beta$ en este problema es no trivial y depende del caso.

%% ====================================================================
\section{Resumen}

\begin{enumerate}
  \item $\alpha$ (solver hibrido): valor intermedio $\sim 0.5$ funciona bien.
  \item En los ejemplos de Nico ($n=3$), aumentado $\approx$ clasico.
        El espacio es demasiado chico para que el conteo agregue valor.
  \item La ventaja del aumentado se manifiesta mejor a $n$ y $B$ mas grandes
        (ver paper\_findings, Parts 2--4).
  \item Greedy secuencial si pierde calidad vs optimo, incluso a $n=3$.
  \item El problema de la ``persona segura'' existe: el efecto de $\beta$
        es no trivial y depende del caso (ver documento 02 para detalles de $\beta$).
\end{enumerate}

\end{document}
